% TOP_PROBLEM: {"id": 6500005, "title": "Symmetries of Vector Distributions", "category_id": 6, "category": {"id": 6, "name": "geometry", "display_name": "Geometry", "description": "Euclidean and non-Euclidean geometry, geometric structures.", "slug": "geometry", "order_index": 6, "created_at": "2026-07-31T15:26:25.671Z"}, "status": "open"}
% FIRST_POSTED: null
% ENTRY_KIND: statement_only
% Imported from: batch13.tex; UnsolvedMath ID 6500005
\documentclass[11pt]{article}
\usepackage[margin=1in]{geometry}
\usepackage{amsmath,amssymb,mathtools}
\usepackage{fontspec,unicode-math}
\setmainfont{Latin Modern Roman}
\setsansfont{Latin Modern Sans}
\setmathfont{Latin Modern Math}
\usepackage{xurl,hyperref}
\hypersetup{colorlinks=true,urlcolor=blue,linkcolor=blue}
\newcommand{\sref}[2]{\hyperlink{source:#1:#2}{[S#2]}}
\newcommand{\cataloguescope}[1]{\paragraph{Recorded mathematical question.}#1}
\begin{document}
% UnsolvedMath ID 6500005; source catalogue code AMR-064-0005
\setcounter{section}{6500004}
\section{Symmetries of Vector Distributions}\label{6500005}
{\small\sffamily UnsolvedMath ID 6500005 · Geometry}
\subsection{Definitions and mathematical statement}

\paragraph{Shared notation.}$\mathbb N=\{1,2,\ldots\}$, and $\mathbb Z,\mathbb Q,\mathbb R,\mathbb C$ denote the integers, rationals, reals and complex numbers. Any other conventions are specified locally.


Let $D\subset TM$ be a smooth vector distribution. Singularity of a horizontal curve means failure of surjectivity of the endpoint-map differential on the space of horizontal paths with fixed starting point. The distribution is singular transitive if every pair of points can be joined by a finite concatenation of singular horizontal curves. A symmetry is a smooth diffeomorphism $\Phi:M\to M$ with $\Phi_*D=D$; no metric is required to be preserved.
\paragraph{Statement recorded in the supplied source.}
Must the full symmetry group of a singular-transitive distribution be a finite-dimensional Lie group? \sref{6500005}{2}
\subsection{Short English statement}
Does connectivity by abnormal curves force the group of distribution-preserving diffeomorphisms to have finite dimension?
\subsection{Sources}
\begin{description}
\item[\hypertarget{source:6500005:1}{S1}] ULAM.AI and the UnsolvedMath Contributors, \emph{UnsolvedMath: A Curated Collection of Open Mathematics Problems}, record 6500005 (AMR-064-0005). CC BY 4.0. \url{https://huggingface.co/datasets/ulamai/UnsolvedMath}.
\item[\hypertarget{source:6500005:2}{S2}] A. A. Agrachev, \emph{Some open problems}, arXiv:1304.2590v2 (2013), Section V, p. 6. \url{https://arxiv.org/pdf/1304.2590}.
\end{description}
\label{end:6500005}
\end{document}
